% Fragmento público generado desde propietarios íntegros.
% Residencia: 52.13.2.
% Fuente: ../incoming/radion_monografia_20260827/manuscrito/sections/13_sintesis.tex:135-184
\subsubsection{System of component theorems of the radion}

\begin{theorem}[Genealogy]
The object \(\mathfrak R_\sigma\) descends from
\(\APP\to\TRIT\to\TPK\to X_{\mathrm{enr}}\) and does not receive the metrological target
as a selector.
\end{theorem}

\begin{theorem}[Hinge]
In the hundred-mark APP chart, the critical seam and the second
dodecaphasic phase satisfy \(j_{100}(1/2)=\theta_2=50^\circ\), with unique \(m=2\).
\end{theorem}

\begin{theorem}[Density]
The two sheets of the active sector \(N_6=90\), normalized by \(3^6=729\),
produce \(\rho_{\mathrm{tw}}=20/81\).
\end{theorem}

\begin{theorem}[Carry]
The APP subtraction of the sections \(S_T\) and \(S_E\) converges to
\(\epsone\) at interval rate \(3^{-54}\) per return and exactly closes
the unit.
\end{theorem}

\begin{theorem}[Action]
The positive ellipse satisfies \(\mathcal A(\theta)=\hret\theta\); a radion
sweeps \(\hret\), and one turn encloses \(\hfull\).
\end{theorem}

\begin{theorem}[Boundary--interior]
The same boundary has finite action \(\hfull\), whereas the Klein
interior has infinite hyperbolic area.
\end{theorem}

\begin{theorem}[Form]
The quotient \(e^{-\eta}\) is preserved across spectral semiaxes,
action semiaxes, uncertainties, impedance, and the modulus of the torus.
\end{theorem}

\begin{theorem}[Polarity]
The inversion \(C^*\mapsto-C^*\) interchanges constitutive relations, inverts
impedance, and preserves the normalized velocity.
\end{theorem}

\begin{theorem}[Modularity]
The bidirectional inversion acts as \(S:\tau\mapsto-1/\tau\) and preserves
\(j(\tau)\), while the focal labeling preserves the orientation that
\(j\) forgets.
\end{theorem}

% Fuente: ../incoming/radion_monografia_20260827/manuscrito/sections/A_demostraciones.tex:1-179
\subsubsection{Algebraic and Geometric Proofs United}

\paragraph{Linear proof of the unit closure}

To facilitate an independent audit, we present the argument without narrative interruption.

\begin{enumerate}[label=\textbf{Step \arabic*.},leftmargin=2.2cm]
\item The nonadic connection generates in a correlated manner
\(\pih,e_{\HMT},\varphi_{\HMT},\alphah\).
\item The turn is \(T_+=2\pih\).
\item The wheel fixes \(\theta_2=50^\circ\), and the parangular chart fixes
\(A=1000\alphah\).
\item The direct section is
\[
S_E=T_+\left(\frac{50}{360}+\frac{1000\alphah}{360}\right)
=2\pih\left(\frac5{36}+\frac{25}{9}\alphah\right).
\]
\item The certified intervals yield \(1<S_E<2\); then
\(D_E=S_E-1\).
\item The incidence produces \(N_6=90\); two sheets and the ambient space \(729\)
produce \(\rho=2N_6/729=20/81\).
\item Global torsion is
\[
\Delta_4=\frac{\pih-e_{\HMT}\log\pih}{270}
\left(1+\frac{\pih}{729}\right).
\]
\item The section torsional is \(S_T=1+(20/81)\Delta_4\).
\item The subtraction APP defines
\[
\epsone=S_T-S_E=\frac{20}{81}\Delta_4-D_E.
\]
\item By substitution,
\[
S_E-\frac{20}{81}\Delta_4+\epsone=1.
\]
\item Multiplying by \(360/T_+=180/\pih\) produces
\[
\frac{180^\circ}{\pih}
=50^\circ+A^\circ-\frac{20}{81}\Delta_4^\circ+\epsR^\circ.
\]
\end{enumerate}

No step receives the left-hand side of the final equality. The
metrological result is evaluated only at the end as a recognition.

\paragraph{Uniqueness of the carry-subtraction}

Let \(B=1000\). For each position, every integer \(u_i\) admits a unique
Euclidean division
\[
u_i=Bc_i+r_i,\qquad 0\le r_i<B.
\]
Therefore,
\[
r_i=u_i\bmod B,\qquad c_i=(u_i-r_i)/B
\]
are unique. Traversing the positions from the remote end toward the
front, the output carry of one position fixes the input of the next.
By induction, the entire word of residues and carries is unique.

In pseudocode:
\begin{verbatim}
acarreo = remote_acarreo
for i from last_block downto first_block:
    u = torsion[i] - direct[i] + acarreo
    remainder[i] = u mod 1000
    acarreo = (u - remainder[i]) / 1000
return remainder, acarreo_word
\end{verbatim}
The procedure does not consult the expected value of the difference.

\paragraph{Convergence of the depth functional}

Let
\[
F(p,e)=1+\frac{20}{81}
\frac{p-e\log p}{270}\left(1+\frac{p}{729}\right)
-2p\left(\frac5{36}+\frac{25}{9}\alphah\right).
\]
On the compact rectangle \(3\le p\le4\), \(2\le e\le3\), \(F\) is
continuous and continuously differentiable. If \(I_N^\pi\searrow\{\pih\}\) and
\(I_N^e\searrow\{e_{\HMT}\}\), the interval image
\[
F(I_N^\pi,I_N^e)
\]
forms a nested family whose width tends to zero. The intersection is a
single real number, \(\epsone\). Monotonicity of the derivatives permits the
endpoints to be evaluated without interior sampling.

The rate \(3^{-54}\) per return follows from the refinement of six trits over
nine levels. At depth \(45\) triads, the width bound
\(<4.81\times10^{-130}\) certifies more than 129 digits of the value.

\paragraph{Derivation of the metric and hyperbolic area of Klein}

On the unit disk, the matrix form of the Klein metric is
\[
g(u,v)=\frac1{1-r^2}I
+\frac1{(1-r^2)^2}
\begin{pmatrix}u\\v\end{pmatrix}
\begin{pmatrix}u&v\end{pmatrix}.
\]
The lemma for the determinant of a rank one perturbation gives
\[
\det g=(1-r^2)^{-3}.
\]
Henceforth
\[
\sqrt{\det g}\,\diff u\diff v
=(1-r^2)^{-3/2}\diff u\diff v.
\]
Integration in polar coordinates yields \eqref{eq:klein-area-R}. The infinite limit is
due to the metric singularity at the boundary; the boundary continues to have
finite symplectic area because that measure uses \(\diff Q\wedge\diff P\),
not \(\sqrt{\det g}\).

\paragraph{Focal Distance on the Klein Disk}

In a diametro, the distance from the center to \(r\) is
\[
d_K(0,r)=\operatorname{artanh}r.
\]
For \(r=e_f\), \eqref{eq:excentricidad} implies
\[
1-e_f^2=e^{-2\eta}.
\]
If \(u_f=\operatorname{artanh}e_f\), then
\[
\operatorname{sech}u_f=\sqrt{1-e_f^2}=e^{-\eta},
\]
and by therefore \(\eta=\log\cosh u_f\), as in
\eqref{eq:focal-hyperbolic-chain}.

\paragraph{Symplectic Compression Transformation}

Define
\begin{equation}
S_\eta=
\begin{pmatrix}e^{\eta/2}&0\\0&e^{-\eta/2}\end{pmatrix},
\qquad \det S_\eta=1.
\label{eq:compresión simpléctica}
\end{equation}
The transported complex structure is
\begin{equation}
J_\eta=S_\eta
\begin{pmatrix}0&-1\\1&0\end{pmatrix}
S_\eta^{-1}
=\begin{pmatrix}0&-e^\eta\\e^{-\eta}&0\end{pmatrix},
\qquad J_\eta^2=-I.
\label{eq:Jeta}
\end{equation}
The curve
\[
\gamma(\theta)=\sqrt{2\hret}\,
S_\eta(\cos\theta,\sin\theta)
=e^{\theta J_\eta}\gamma(0)
\]
is exactly the action ellipse. The hyperbolic symplectic compression preserves area and
transports complex rotation instead of destroying it.

For orientation \(\sigma\),
\[
\gamma_\sigma(\theta)
=(a_S\cos\theta,\sigma b_S\sin\theta),
\qquad
\mathcal A_\sigma(\theta)=\sigma\hret\theta.
\]
The two sheets share foci and action modulus; they differ in signed action.

\paragraph{Diophantine verification of the dodecaphasic Barbero--Immirzi weight}

The equation \(4a-3b=45\) has integer solutions
\[
(a,b)=(12+3t,1+4t),\qquad t\in\Z.
\]
The condition \(|b|=1\) leaves \(b=1\), since \(b=-1\) does not belong to the class
\(1\pmod4\). Hence \(a=12\). This verification reproduces the pair obtained through the dodecaphasic
fundamental chain: twelve sectors and one unique oriented seam. The cardinalities
\(90/120\), the chain multiplicities, and the return reader
retain distinct functions; the pole coefficient and minimality
arithmetically verify their composition.
